\documentclass[11pt,a4paper]{article}

% ── Packages ──────────────────────────────────────────────────────────
\usepackage[utf8]{inputenc}
\usepackage[T1]{fontenc}
\usepackage{amsmath,amssymb}
\usepackage{booktabs}
\usepackage{graphicx}
\usepackage{xcolor}
\usepackage{hyperref}
\usepackage{enumitem}
\usepackage{fancyhdr}
\usepackage{geometry}
\usepackage{titlesec}
\usepackage{mdframed}
\usepackage{listings}
\usepackage{float}

\geometry{margin=1in}
\hypersetup{colorlinks=true,linkcolor=blue!60!black,urlcolor=blue!60!black}

% ── Colors ────────────────────────────────────────────────────────────
\definecolor{accent}{HTML}{1E3A8A}
\definecolor{lightgray}{HTML}{F3F4F6}
\definecolor{codebg}{HTML}{F8FAFC}

% ── Title formatting ──────────────────────────────────────────────────
\titleformat{\section}{\Large\bfseries\color{accent}}{}{0em}{}
\titleformat{\subsection}{\large\bfseries\color{accent!80!black}}{}{0em}{}
\titleformat{\subsubsection}{\normalsize\bfseries\color{accent!60!black}}{}{0em}{}

% ── Code block ────────────────────────────────────────────────────────
\lstset{
  basicstyle=\ttfamily\small,
  backgroundcolor=\color{codebg},
  frame=single,
  framerule=0.5pt,
  rulecolor=\color{lightgray},
  breaklines=true,
  xleftmargin=1em,
  xrightmargin=1em,
  aboveskip=1em,
  belowskip=1em,
}

% ── Callout box ───────────────────────────────────────────────────────
\newmdenv[
  linecolor=accent!40,
  backgroundcolor=accent!5,
  roundcorner=5pt,
  skipabove=1em,
  skipbelow=1em,
  innerleftmargin=1em,
  innerrightmargin=1em
]{callout}

% ── Header / Footer ───────────────────────────────────────────────────
\pagestyle{fancy}
\fancyhf{}
\fancyhead[L]{\small\color{gray}Rare Disease Pre-diagnosis Clue Finder · Architecture Document}
\fancyhead[R]{\small\color{gray}iGEM 2026 PekingHSC}
\fancyfoot[C]{\thepage}
\renewcommand{\headrulewidth}{0.4pt}

\begin{document}

% ═══════════════════════════════════════════════════════════════════════
% TITLE PAGE
% ═══════════════════════════════════════════════════════════════════════
\begin{titlepage}
\vspace*{3cm}
\begin{center}
{\Huge\bfseries\color{accent} Rare Disease Pre-diagnosis Clue Finder}\\[0.5cm]
{\Large Architecture \& Algorithm Reference}\\[1.5cm]
{\large iGEM 2026 · PekingHSC}\\[0.3cm]
{\large Human Phenotype Ontology × Orphanet Knowledge Graph}\\[2cm]
{\normalsize Wei Ziheng (Sebastian Wei)}\\[0.3cm]
{\small \today}
\end{center}
\vfill
\begin{callout}
\textbf{Abstract.} This document describes the complete architecture of the Rare Disease Pre-diagnosis Clue Finder—a knowledge-graph-powered diagnostic assistant that translates colloquial symptom descriptions into structured HPO (Human Phenotype Ontology) terms, retrieves candidate matches via a hybrid keyword–vector index, validates them through an LLM-based reranking stage, and infers the most likely rare diseases using a True Path Rule engine with IDF-weighted scoring. The system is designed to run on a 2-core, 4\,GB server with all NLP and embedding workloads offloaded to cloud APIs.
\end{callout}
\end{titlepage}

% ═══════════════════════════════════════════════════════════════════════
\tableofcontents
\newpage

% ═══════════════════════════════════════════════════════════════════════
\section{System Overview}
% ═══════════════════════════════════════════════════════════════════════

\subsection{What Problem Does It Solve?}

A patient walks into a clinic and says:

\begin{quote}
\itshape
``My hands and feet feel like they're burning. I have these dark red spots on my belly and thighs. I can't sweat at all. Every time I eat, I get diarrhea. I always feel overheated.''
\end{quote}

A general practitioner might order routine tests. But these seven words—\textbf{burning pain}, \textbf{dark red spots}, \textbf{inability to sweat}, \textbf{postprandial diarrhea}, \textbf{heat intolerance}—are, together, the textbook presentation of \textbf{Fabry disease}, a rare lysosomal storage disorder with a median diagnostic delay of over a decade.

The Rare Disease Pre-diagnosis Clue Finder automates exactly this pattern-recognition process. It takes free-text, colloquial symptom descriptions and maps them through a structured medical ontology to produce a ranked list of candidate diseases, each accompanied by an \textit{explainable evidence trail}.

\subsection{The Pipeline at a Glance}

\begin{enumerate}[leftmargin=*,itemsep=4pt]
  \item \textbf{Symptom Extraction.} A large language model (DeepSeek v4-flash) reads the free-text input, breaks it into atomic symptoms, and translates each into a standardized HPO term.
  \item \textbf{Candidate Retrieval.} Each standardized term is looked up in two indices simultaneously: a lexical keyword index (searching HPO names, synonyms, and definitions) and a FAISS vector index (1024-dimensional embeddings via BAAI/bge-m3). Results are merged into a unified candidate pool of up to 25 terms.
  \item \textbf{LLM Reranking.} DeepSeek reviews each candidate pool \textit{with the full original patient text as context}, selects the single best HPO match for each symptom, or rejects all candidates if none fit.
  \item \textbf{True Path Rule Inference.} The selected HPO terms are fed into a NetworkX knowledge graph (35,964 nodes, 139,516 edges). The graph traverses the HPO hierarchy—accounting for direct, child, and parent matches—and scores each candidate disease using frequency-weighted contributions with an IDF-smoothing penalty that prevents common symptoms from dominating.
  \item \textbf{Explainable Output.} The top 5 diseases are returned, each with a full evidence trail showing \textit{which} symptom matched \textit{which} HPO term, \textit{how} (direct/child/parent), and \textit{how much} it contributed to the final score.
\end{enumerate}

\subsection{Technology Stack}

\begin{table}[H]
\centering
\begin{tabular}{@{}lll@{}}
\toprule
\textbf{Layer} & \textbf{Technology} & \textbf{Role} \\
\midrule
Inference Engine & NetworkX DiGraph & Directed knowledge graph \\
Symptom Extraction & DeepSeek v4-flash API & Zero-shot medical NER + HPO standardization \\
Vector Retrieval & FAISS IndexFlatIP + BAAI/bge-m3 & 1024-dim semantic search \\
Candidate Reranking & DeepSeek v4-flash API & Context-aware medical candidate selection \\
Web Framework & FastAPI + uvicorn & Async REST API \\
Data Sources & HPO v2026-02-16 + Orphanet ORDO v4.8 & Phenotype ontology + rare disease KB \\
\bottomrule
\end{tabular}
\end{table}

\newpage

% ═══════════════════════════════════════════════════════════════════════
\section{The Knowledge Graph}
% ═══════════════════════════════════════════════════════════════════════

\subsection{Graph Topology}

The knowledge graph is a \textbf{heterogeneous directed graph} with two node types and two edge types:

\begin{table}[H]
\centering
\begin{tabular}{@{}llll@{}}
\toprule
\textbf{Type} & \textbf{Count} & \textbf{ID Format} & \textbf{Attributes} \\
\midrule
Phenotype nodes & 19,944 & \texttt{HP:NNNNNNN} & name, synonyms, definition \\
Disease nodes   & 16,020 & \texttt{ORPHA:NNNN} & name \\
\midrule
\textit{Total nodes} & \textit{35,964} & & \\
\bottomrule
\end{tabular}
\end{table}

\begin{table}[H]
\centering
\begin{tabular}{@{}llll@{}}
\toprule
\textbf{Type} & \textbf{Count} & \textbf{Direction} & \textbf{Meaning} \\
\midrule
\texttt{is\_a} & 23,677 & child $\rightarrow$ parent & HPO ontological hierarchy \\
\texttt{disease\_has\_hpo} & 115,878 & disease $\rightarrow$ hpo & Disease–phenotype association (with frequency weight) \\
\midrule
\textit{Total edges} & \textit{139,516} & & \\
\bottomrule
\end{tabular}
\end{table}

\begin{callout}
\textbf{Direction convention.} HPO \texttt{is\_a} edges point from child (specific) to parent (general). This means \texttt{nx.ancestors(G, hpo)} returns \textit{children} (more specific terms), and \texttt{nx.descendants(G, hpo)} returns \textit{parents} (more general terms). The code uses semantic naming—``children'' and ``parents'' refer to the clinical relationship, not the graph direction.
\end{callout}

\subsection{Frequency Weights}

Each disease–HPO edge carries a frequency weight indicating how often the phenotype appears in that disease:

\begin{table}[H]
\centering
\begin{tabular}{@{}ll@{}}
\toprule
\textbf{Frequency Label} & \textbf{Weight} \\
\midrule
Obligate (100\%)     & 1.00 \\
Very frequent (80–99\%) & 0.90 \\
Frequent (30–79\%)    & 0.55 \\
Occasional (5–29\%)   & 0.17 \\
Very rare ($<$5\%)     & 0.02 \\
Excluded (0\%)        & 0.00 \\
No data               & 0.50 \\
\bottomrule
\end{tabular}
\end{table}

\subsection{Data Ingestion}

The script \texttt{parse\_data.py} processes three raw sources into three CSV files:

\begin{itemize}[leftmargin=*]
  \item \texttt{hp.json} — HPO ontology: term names, synonyms, \textbf{definitions}, and \texttt{is\_a} edges.
  \item \texttt{ORDO\_en\_4.8.owl.xml} — Orphanet Rare Disease Ontology: disease names and IDs.
  \item \texttt{en\_product4.xml} — Orphanet epidemiological data: disease–HPO associations with frequency labels.
\end{itemize}

The output \texttt{nodes.csv} includes a \texttt{definition} column containing the full HPO definition text for each phenotype term. This field is critical for downstream retrieval: for example, ``Acroparesthesia'' has an empty synonym list but its definition reads \textit{``burning or numbness\ldots in the hands and feet''}—which directly bridges the layperson phrase ``burning pain in my hands.''

\newpage

% ═══════════════════════════════════════════════════════════════════════
\section{Symptom Extraction \& Standardization}
% ═══════════════════════════════════════════════════════════════════════

\subsection{The Semantic Gap Problem}

Patients describe symptoms in colloquial language. The HPO ontology uses precise medical nomenclature. The gap between ``can't sweat'' and ``Anhidrosis''—or between ``dark red spots'' and ``Angiokeratoma''—is too wide for simple string matching or naive vector search to bridge reliably.

\subsection{DeepSeek Extraction Prompt}

We use a single zero-shot prompt to perform both entity extraction and HPO standardization:

\begin{lstlisting}[caption={Extraction prompt (simplified)},label={lst:extract}]
You are an expert clinical information extractor.
Analyze the following user text and extract the distinct
medical signs and symptoms.

Rules:
1. Break down complex descriptions into atomic symptoms.
2. MUST include anatomical location.
3. For 'standard_term', output the EXACT HPO term name
   whenever possible — not a general clinical description.
   Examples:
     "dark red spots on skin" → "Angiokeratoma"
                                (not "purpura" or "red spots")
     "burning pain in hands and feet" → "Acroparesthesia"
                                (not "burning paresthesia")
     "cannot sweat" → "Anhidrosis"
                                (not "absence of sweating")
4. Return a JSON array of {raw_description, standard_term}.
\end{lstlisting}

\subsection{Output Structure}

\begin{lstlisting}[caption={Example extraction output}]
[
  {"raw_description": "burning pain in hands",
   "standard_term":   "Acroparesthesia"},
  {"raw_description": "dark red spots on belly",
   "standard_term":   "Angiokeratoma"},
  {"raw_description": "cannot sweat at all",
   "standard_term":   "Anhidrosis"}
]
\end{lstlisting}

\subsection{The Critical Design Decision: Strict Separation}

The two fields serve different consumers and must \textbf{never be concatenated}:

\begin{itemize}[leftmargin=*]
  \item \texttt{raw\_description} is for the \textbf{user interface}. The patient sees their own words mirrored back, building trust.
  \item \texttt{standard\_term} is for the \textbf{retrieval pipeline}. It feeds the keyword index, the FAISS index, and the DeepSeek reranker.
\end{itemize}

Concatenating them (e.g., \texttt{"burning pain in hands (Acroparesthesia)"}) causes \textbf{semantic dilution}: the embedding model's attention is scattered across high-frequency layperson words (\textit{hands}, \textit{feet}, \textit{pain}), drowning out the precise medical signal needed for HPO ontology matching. This single architectural decision eliminated the majority of false-positive vector matches observed in early testing.

\newpage

% ═══════════════════════════════════════════════════════════════════════
\section{The Retrieval System}
% ═══════════════════════════════════════════════════════════════════════

\subsection{Dual-Index Architecture}

Each standardized symptom term is looked up in two independent indices, and the results are merged into a single candidate pool:

\begin{center}
\begin{tabular}{ccc}
\textbf{Keyword Index} & \qquad & \textbf{FAISS Vector Index} \\
\small(graph.search\_hpo) & & \small(IndexFlatIP, 1024-dim) \\
$\downarrow$ & & $\downarrow$ \\
\small Substring match against & & \small Top-20 nearest neighbors \\
\small name, synonyms, \textbf{definition} & & \small via BAAI/bge-m3 embedding \\
$\downarrow$ & & $\downarrow$ \\
\multicolumn{3}{c}{\textbf{Unified Candidate Pool ($\leq$ 25 terms)}} \\
\end{tabular}
\end{center}

\subsection{Keyword Index}

The keyword index performs case-insensitive substring matching with a three-tier fallback:

\begin{enumerate}[leftmargin=*]
  \item \textbf{Name.} Exact or partial match against the HPO term name (e.g., ``Anhidrosis'' $\leftarrow$ ``anhidrosis'').
  \item \textbf{Synonyms.} If the name fails, search the synonyms field (e.g., ``Lack of sweating'' $\leftarrow$ ``cannot sweat'').
  \item \textbf{Definition.} If synonyms also fail, search the definition text (e.g., ``burning\ldots in the hands and feet'' $\leftarrow$ ``burning pain in hands'').
\end{enumerate}

The definition search is what makes the keyword index genuinely useful. Without it, ``Acroparesthesia'' (which has no synonyms in the ontology) would be invisible to keyword search for queries like ``burning pain''—even though its definition explicitly contains those words.

\subsection{FAISS Vector Index}

\begin{table}[H]
\centering
\begin{tabular}{@{}ll@{}}
\toprule
\textbf{Parameter} & \textbf{Value} \\
\midrule
Embedding model & BAAI/bge-m3 (SiliconFlow API) \\
Dimensionality & 1024 \\
Index type & IndexFlatIP (inner product = cosine on L2-normalized vectors) \\
Index size & 19,944 HPO phenotype terms \\
Query side & With instruction prefix \\
Document side & Without prefix \\
\bottomrule
\end{tabular}
\end{table}

\subsubsection{Document Text Format}

Each HPO term in the FAISS index is embedded as a structured text block:

\begin{lstlisting}[caption={Index text for Acroparesthesia (HP:0031006)}]
Name: Acroparesthesia | Definition: A type of paresthesia
(tingling, pins-and-needles, burning or numbness or stiffness)
that occurs in the hands and feet and particularly in the
fingers and toes.
\end{lstlisting}

Embedding only the term name (``Acroparesthesia'') would give the vector model no signal that this term is related to ``burning'' or ``hands.'' Including the definition bridges this gap: the embedding now encodes both the medical label \textit{and} the layperson-accessible description.

\subsubsection{Asymmetric Retrieval Design}

The BGE model family supports instruction-tuned encoding, where the query and document sides receive different prompts:

\begin{itemize}[leftmargin=*]
  \item \textbf{Query side} (what we search \textit{with}):\\
  \texttt{"Represent this medical symptom for retrieving relevant rare disease}\\\texttt{phenotype terms from a standard ontology: Acroparesthesia"}
  \item \textbf{Document side} (what we search \textit{in}):\\
  \texttt{"Name: Acroparesthesia | Definition: \ldots"} \hfill (\textit{no prefix})
\end{itemize}

This asymmetry signals the model to encode the query as a \textit{search intent} and the document as a \textit{passive passage}, improving recall for short queries against longer document texts.

\newpage

% ═══════════════════════════════════════════════════════════════════════
\section{Two-Stage Candidate Reranking}
% ═══════════════════════════════════════════════════════════════════════

\subsection{Why Reranking Is Necessary}

FAISS is excellent at recall but terrible at precision for medical text. It operates purely on vector similarity—it has no understanding of medicine. Given ``hand pain'' as a query, FAISS might recall ``Fused finger joints'' because both involve \textit{hand} and \textit{pain} in similar vector neighborhoods. A human clinician (or a capable LLM) instantly knows this is wrong.

Conversely, an LLM asked to suggest an HPO term from scratch might \textit{hallucinate}—inventing a term that does not exist in the ontology.

The two-stage architecture cancels out both failure modes:

\begin{callout}
\textbf{FAISS constrains the LLM.} The LLM can only choose from real HPO terms that actually exist in the ontology—eliminating hallucinations.\\
\textbf{The LLM corrects FAISS.} The LLM reads the clinical context and rejects vector-neighbor false positives—eliminating nonsensical matches.
\end{callout}

\subsection{Unified Candidate Pool (v3 Architecture)}

Earlier versions used a ``keyword fast-path'' where keyword matches were accepted directly without LLM review. This was removed in the current architecture. The v3 flow collects all candidates—keyword and FAISS—into a single pool and sends everything through the same DeepSeek gatekeeper:

\begin{lstlisting}[caption={Unified candidate collection (per symptom)}]
# Step 1: Collect candidates from both sources
candidates = []

# a. Keyword search (limit=5, searches name/synonyms/definition)
kw_hits = graph.search_hpo(standard_term, limit=5)
candidates.extend(kw_hits)

# b. FAISS vector recall (top-20)
if vector_index is not None:
    query_vec = get_api_embedding(standard_term)
    distances, indices = vector_index.search(query_vec, 20)
    candidates.extend(faiss_results)

# c. Deduplicate
candidates = deduplicate(candidates)  # ≤ 25 unique terms

# Step 2: LLM rerank (single gatekeeper for all candidates)
best_id = llm_rerank_candidates(
    raw_description,      # single symptom fragment
    candidates,            # unified pool
    full_text=original     # full patient narrative as context
)
\end{lstlisting}

\subsection{Reranker Prompt Design}

The reranker receives three pieces of information:

\begin{enumerate}[leftmargin=*]
  \item The single symptom fragment being resolved (e.g., ``dark red spots on belly and thighs'')
  \item The \textbf{full original patient text} as clinical context (e.g., ``burning pain in hands, dark red spots, cannot sweat, diarrhea after eating, feeling overheated'')
  \item The unified candidate pool ($\leq$ 25 HPO terms)
\end{enumerate}

The full text is the key. In isolation, ``dark red spots'' could reasonably match ``Purpura'' (a generic bleeding-under-skin term). But when the model also sees ``burning pain in extremities'' + ``cannot sweat'' + ``postprandial diarrhea,'' the clinical picture shifts decisively toward Fabry disease, and the correct match becomes ``Angiokeratoma.''

\begin{lstlisting}[caption={Reranker prompt (abbreviated)}]
You are an expert medical ontology matcher.
The patient reported: "dark red spots on belly and thighs"

The patient's FULL description was: "burning pain in hands,
dark red spots, cannot sweat, diarrhea after eating, overheated"

Candidates:
- ID: HP:0000979 | Term: Purpura
- ID: HP:0001014 | Term: Angiokeratoma
- ...

TASK: Select the SINGLE candidate that best fits.
Use the FULL clinical context to disambiguate.
If none match, reject all.

Output: {"selected_id": "HP:0001014"} OR {"selected_id": null}
\end{lstlisting}

\subsection{Fault Tolerance}

\begin{table}[H]
\centering
\begin{tabular}{@{}ll@{}}
\toprule
\textbf{Condition} & \textbf{Behavior} \\
\midrule
DeepSeek API succeeds & Selects best HPO or returns \texttt{null} \\
DeepSeek API fails & Falls back to top keyword match (\texttt{keyword\_fallback}) \\
Candidate pool is empty & Marks symptom as \texttt{(no match found)} \\
Symptom unmatched & Excluded from inference; does not affect other symptoms \\
\bottomrule
\end{tabular}
\end{table}

\newpage

% ═══════════════════════════════════════════════════════════════════════
\section{The Inference Engine}
% ═══════════════════════════════════════════════════════════════════════

\subsection{True Path Rule}

The fundamental inference rule of the HPO ontology:

\begin{callout}
\textbf{True Path Rule.} If a disease is associated with a phenotype term $P$, it is implicitly associated with all ancestors (more general terms) and all descendants (more specific terms) of $P$ in the HPO hierarchy.
\end{callout}

For example, if the user selects \texttt{HP:0001250} (Seizure):
\begin{itemize}[leftmargin=*]
  \item \textbf{Direct match:} The disease is directly linked to Seizure $\rightarrow$ full weight.
  \item \textbf{Child match:} The disease is linked to ``Generalized tonic-clonic seizure'' (a child of Seizure) $\rightarrow$ reduced weight.
  \item \textbf{Parent match:} The disease is linked to ``Abnormal nervous system physiology'' (a parent of Seizure) $\rightarrow$ reduced weight.
\end{itemize}

\subsection{Scoring Formula}

For a given set of input HPO terms, each candidate disease receives a cumulative score:

\begin{equation}
\text{Score}(\text{disease}) = \sum_{\text{matched HPO}} \Big[ \text{freq\_weight} \times \text{multiplier} \times \text{idf\_weight} \Big]
\end{equation}

Where the multiplier depends on match type:

\begin{equation}
\text{multiplier} = \begin{cases}
1.0 & \text{direct match} \\[4pt]
0.5 \times e^{-(\text{rank}-1)} & \text{indirect match (child or parent)}
\end{cases}
\end{equation}

The rank-decay mechanism for indirect matches: within each (disease, input HPO) group, indirect matches are sorted by frequency weight descending, then weighted exponentially. Rank 1 gets 0.5, rank 2 gets $0.5 \times e^{-1} \approx 0.184$, rank 3 gets $0.5 \times e^{-2} \approx 0.068$, and so on.

\begin{callout}
\textbf{Why rank-decay?} Without it, HPO terms with large descendant trees (e.g., ``Abnormality of the nervous system'' with $50+$ children) would dominate the ranking. Rank-decay ensures only the top 1–2 best-matching children per disease contribute meaningfully.
\end{callout}

\subsection{IDF Smoothing: Why Standard IDF Fails in Medicine}

Standard TF-IDF uses the raw logarithm:

\begin{equation}
\text{IDF}_{\text{standard}} = \ln\left(\frac{N_{\text{total}}}{df + 1}\right)
\end{equation}

With $N \approx 16{,}000$ diseases in the graph, this produces extreme ratios:

\begin{table}[H]
\centering
\begin{tabular}{@{}lcc@{}}
\toprule
\textbf{Symptom} & \textbf{df (disease count)} & \textbf{Standard IDF} \\
\midrule
Fever (common)       & 3,000 & 1.67 \\
Abdominal pain       & 2,000 & 2.08 \\
Seizure              & 800   & 3.00 \\
Purpura              & 200   & 4.38 \\
Anhidrosis (rare)    & 30    & 6.25 \\
Angiokeratoma (rare) & 15    & 6.90 \\
\bottomrule
\end{tabular}
\end{table}

A single rare symptom (IDF $\approx$ 6.9) contributes nearly \textbf{4 times} as much as a common one (IDF $\approx$ 1.7). If the LLM incorrectly matches a rare term, or if the patient happens to mention an incidental rare finding, that single term can single-handedly dominate the ranking—pushing the correct disease out of the top 5.

\subsection{Square-Root Smoothing}

We apply a compressive transformation:

\begin{equation}
\text{IDF}_{\text{smoothed}} = 1.0 + \sqrt{\;\ln\left(\frac{N_{\text{total}}}{df + 1}\right)\;}
\end{equation}

\begin{table}[H]
\centering
\begin{tabular}{@{}lccc@{}}
\toprule
\textbf{Symptom} & \textbf{Standard IDF} & \textbf{Smoothed IDF} & \textbf{Rare/Common Ratio} \\
\midrule
Fever               & 1.67 & 2.29 & — \\
Abdominal pain      & 2.08 & 2.44 & 1.07$\times$ \\
Seizure             & 3.00 & 2.73 & 1.19$\times$ \\
Purpura             & 4.38 & 3.09 & 1.35$\times$ \\
Anhidrosis          & 6.25 & 3.50 & 1.53$\times$ \\
Angiokeratoma       & 6.90 & 3.63 & 1.58$\times$ \\
\bottomrule
\end{tabular}
\end{table}

The $\sqrt{\cdot}$ function compresses the rare/common ratio from \textbf{4.1$\times$} down to \textbf{1.6$\times$}. Rare symptoms are still rewarded—they \textit{do} pull rare diseases up the ranking—but they can no longer single-handedly override a disease with several moderately specific matches.

\begin{callout}
\textbf{Design rationale.} The $+1.0$ base ensures that \textit{every} match contributes at least one unit of score. The $\sqrt{\log(\cdot)}$ term provides a bounded ``rarity bonus'' that distinguishes discriminative symptoms from generic ones without creating instability. Three medium-specificity symptoms (smoothed IDF $\approx$ 2.5 each, total 7.5) will always beat one ultra-rare symptom (smoothed IDF $\approx$ 3.6).
\end{callout}

\subsection{Dynamic Symptom Elicitation}

After each inference, the system suggests the next most diagnostically useful symptoms to ask about:

\begin{enumerate}[leftmargin=*]
  \item Collect all HPO terms linked to the top-5 candidate diseases.
  \item Exclude terms already covered by the True Path Rule (the user's selected HPOs and their ancestors/descendants are already accounted for).
  \item Score remaining terms by \textbf{discriminative power}:\\
  $\text{score} = \max\_\text{frequency} \times \text{discriminative\_factor}$\\
  where discriminative\_factor = 1.0 (appears in 1 disease), 0.5 (2 diseases), 0.25 (3+ diseases).
  \item Return the top 10.
\end{enumerate}

A symptom that is Obligate (frequency 1.0) in exactly one of the top-5 diseases is the ideal discriminator—asking about it will immediately confirm or rule out that disease.

\newpage

% ═══════════════════════════════════════════════════════════════════════
\section{API Design}
% ═══════════════════════════════════════════════════════════════════════

\subsection{Endpoint Summary}

\begin{table}[H]
\centering
\begin{tabular}{@{}lll@{}}
\toprule
\textbf{Endpoint} & \textbf{Method} & \textbf{Purpose} \\
\midrule
\texttt{/} & GET & Frontend SPA \\
\texttt{/api/health} & GET & Health check + graph statistics \\
\texttt{/api/predict} & POST & Manual HPO selection $\rightarrow$ inference \\
\texttt{/api/hpo-search} & GET & Keyword search (name + synonyms + definition) \\
\texttt{/api/smart-search} & GET & Extraction + hybrid search (guided mode) \\
\texttt{/api/vector-search} & GET & Pure FAISS search (debugging) \\
\texttt{/api/auto-diagnose} & GET & \textbf{Zero-click full pipeline} (default for patients) \\
\texttt{/api/anatomy-tree} & GET & Top-level HPO system categories \\
\bottomrule
\end{tabular}
\end{table}

\subsection{The Auto-Diagnose Endpoint}

This is the primary endpoint for patient-facing use. It accepts free text and returns a complete diagnostic report.

\textbf{Request:}
\begin{lstlisting}
GET /api/auto-diagnose?text=burning+pain+in+hands+dark+red+spots+cannot+sweat&mode=public
\end{lstlisting}

\textbf{Response} (abbreviated):
\begin{lstlisting}
{
  "query_hpo_ids": ["HP:0031006", "HP:0001014", "HP:0000970"],
  "results": [
    {
      "disease_id": "ORPHA:324",
      "disease_name": "Fabry disease",
      "total_score": 8.52,
      "matched_paths": [
        {
          "input_hpo_id": "HP:0031006",
          "match_type": "direct",
          "contribution": 3.15,
          "explanation": "matched directly with Acroparesthesia [IDF: 3.50]"
        }
      ]
    }
  ],
  "auto_selections": [
    {
      "raw_description": "burning pain in hands",
      "standard_term": "Acroparesthesia",
      "matched_hpo_id": "HP:0031006",
      "matched_hpo_name": "Acroparesthesia",
      "match_method": "llm_rerank"
    }
  ]
}
\end{lstlisting}

The \texttt{auto\_selections} array provides full transparency: the user can see exactly which of their words was mapped to which HPO term, and whether that mapping came from the LLM reranker (\texttt{llm\_rerank}) or a keyword fallback (\texttt{keyword\_fallback}).

\subsection{Dual-Mode Architecture}

The system supports two display modes via the \texttt{?mode=} parameter:

\begin{table}[H]
\centering
\begin{tabular}{@{}lll@{}}
\toprule
\textbf{Feature} & \textbf{Expert Mode} & \textbf{Public Mode (default)} \\
\midrule
HPO names & Original medical terms & Layperson translations \\
HPO IDs & Visible (\texttt{HP:XXXXXXX}) & Hidden \\
Cytoscape animation & Shown & Skipped (direct to report) \\
Manual HPO search & Available & Hidden (Smart Match / Auto-Diagnose only) \\
Auto-Diagnose toggle & Hidden & Visible (default: ON) \\
\bottomrule
\end{tabular}
\end{table}

\newpage

% ═══════════════════════════════════════════════════════════════════════
\section{Deployment}
% ═══════════════════════════════════════════════════════════════════════

\subsection{Hardware Requirements}

\begin{table}[H]
\centering
\begin{tabular}{@{}ll@{}}
\toprule
\textbf{Resource} & \textbf{Minimum} \\
\midrule
CPU & 2 cores \\
RAM & 4 GB (graph $\approx$ 140 MB + FAISS index $\approx$ 80 MB) \\
Disk & $\approx$ 500 MB (data files + index) \\
\bottomrule
\end{tabular}
\end{table}

All NLP and embedding computation is offloaded to cloud APIs (DeepSeek + SiliconFlow). No local models are loaded. The server starts in under 2 seconds.

\subsection{Environment Variables}

\begin{table}[H]
\centering
\begin{tabular}{@{}ll@{}}
\toprule
\textbf{Variable} & \textbf{Purpose} \\
\midrule
\texttt{DEEPSEEK\_API\_KEY} & DeepSeek API key (symptom extraction + reranking) \\
\texttt{EMBEDDING\_API\_KEY} & SiliconFlow API key (BAAI/bge-m3 embeddings) \\
\texttt{EMBEDDING\_BASE\_URL} & Embedding API base URL (default: \texttt{api.siliconflow.cn/v1}) \\
\texttt{EMBEDDING\_MODEL} & Embedding model name (default: \texttt{BAAI/bge-m3}) \\
\bottomrule
\end{tabular}
\end{table}

\subsection{Startup Sequence}

\begin{lstlisting}[caption={Deployment commands}]
cd /path/to/RD

# Install dependencies
pip install -r requirements.txt

# Set API keys
export DEEPSEEK_API_KEY="sk-..."
export EMBEDDING_API_KEY="sk-..."

# Build FAISS index (first time, or after nodes.csv changes)
python3 build_vector_index.py

# Generate layperson translations (optional, for public mode)
python3 build_layperson_dict.py

# Start server
python3 main.py
# or: python3 -m uvicorn main:app --host 0.0.0.0 --port 8000
\end{lstlisting}

\subsection{Concurrency Considerations}

Each \texttt{/api/auto-diagnose} request makes $N+1$ serial API calls to DeepSeek (1 extraction + $N$ reranks, where $N$ is the number of extracted symptoms). With a single uvicorn worker, one request occupies the worker for 3–10 seconds depending on network latency. For production multi-user deployment:

\begin{itemize}[leftmargin=*]
  \item Use multiple workers: \texttt{uvicorn main:app --workers 4}
  \item Wrap synchronous API calls in \texttt{asyncio.to\_thread()} to avoid blocking the event loop.
  \item Parallelize the $N$ reranker calls with \texttt{asyncio.gather()}.
\end{itemize}

\newpage

% ═══════════════════════════════════════════════════════════════════════
\section{File Inventory}
% ═══════════════════════════════════════════════════════════════════════

\begin{table}[H]
\centering
\begin{tabular}{@{}ll@{}}
\toprule
\textbf{File} & \textbf{Role} \\
\midrule
\texttt{parse\_data.py} & Raw ontology files $\rightarrow$ CSV graph data \\
\texttt{main.py} & FastAPI backend: graph, endpoints, LLM functions \\
\texttt{index.html} & Single-file frontend (dossier-style, auto-diagnose toggle) \\
\texttt{build\_vector\_index.py} & Build FAISS index via cloud embedding API \\
\texttt{build\_layperson\_dict.py} & Generate layperson translation dictionary \\
\texttt{nodes.csv} & 35,964 nodes (includes \texttt{definition} column) \\
\texttt{edges\_hpo.csv} & 23,677 HPO hierarchy edges \\
\texttt{edges\_disease\_hpo.csv} & 115,878 disease–HPO association edges \\
\texttt{hpo\_index.faiss} & FAISS index (1024-dim, $\approx$ 80 MB) \\
\texttt{hpo\_mapping.pkl} & Index position $\rightarrow$ HPO ID/name mapping \\
\texttt{layperson\_dict.json} & HPO $\rightarrow$ layperson phrase translations \\
\texttt{requirements.txt} & Python dependencies \\
\bottomrule
\end{tabular}
\end{table}

\newpage

% ═══════════════════════════════════════════════════════════════════════
\section{Design Decision Record}
% ═══════════════════════════════════════════════════════════════════════

\begin{table}[H]
\centering
\small
\begin{tabular}{@{}p{3.8cm}p{4.2cm}p{4.5cm}@{}}
\toprule
\textbf{Decision} & \textbf{Approach} & \textbf{Rationale} \\
\midrule
Raw/standard separation & Two independent fields; never concatenated for embedding & Semantic dilution: layperson high-frequency words scatter model attention \\
\midrule
Keyword search includes definition & \texttt{search\_hpo} matches against name, synonyms, \textit{and} definition & HPO definitions contain layperson-accessible descriptions (e.g., ``burning\ldots in the hands'') \\
\midrule
No keyword fast-path & Keyword and FAISS results merged; all go through DeepSeek reranker & Keyword is too broad post-definition-search; LLM review prevents false positives \\
\midrule
Reranker receives full patient text & \texttt{llm\_rerank\_candidates(full\_text=original)} & Single symptom ambiguous in isolation; full clinical picture enables correct disambiguation \\
\midrule
IDF sqrt smoothing & $\text{IDF} = 1 + \sqrt{\ln(N/df)}$ & Standard log-IDF is too steep for medicine (rare symptom = 4$\times$ common); sqrt compresses to 1.6$\times$ \\
\midrule
Extraction prompt with counter-examples & Show \textit{what not to output}: ``Angiokeratoma (not purpura)'' & LLMs default to generic clinical terms; explicit negatives steer toward precise HPO names \\
\midrule
FAISS index includes definition & Embed \texttt{Name: X | Definition: ...} not just name & Definition contains layperson synonyms; bridges the semantic gap that pure-name embedding cannot \\
\midrule
Instruction prefix on query only & Prefix added to query embeddings; index embeddings have no prefix & Asymmetric BGE encoding: query-as-search-intent vs. document-as-passive-passage \\
\bottomrule
\end{tabular}
\end{table}

\vspace{1cm}
\begin{center}
\rule{0.4\textwidth}{0.4pt}\\[0.5cm]
{\small iGEM 2026 · PekingHSC · CC-BY-4.0}
\end{center}

\end{document}
